% ============================================================================
% English UI and structural strings for the One Computer Science Book series.
% Loaded by styles/onecs.sty when \booklang is en (default).
% ============================================================================

% Theorem / statement environment titles
\newcommand{\omnameDefinition}{Definition}
\newcommand{\omnameTheorem}{Theorem}
\newcommand{\omnameProposition}{Proposition}
\newcommand{\omnameLemma}{Lemma}
\newcommand{\omnameCorollary}{Corollary}
\newcommand{\omnameMethod}{Method}
\newcommand{\omnameExample}{Example}
\newcommand{\omnameNotation}{Notation}
\newcommand{\omnameRemark}{Remark}
\newcommand{\omnameExercise}{Exercise}
\newcommand{\omnameExercises}{Exercises}
\newcommand{\omnameProblem}{Problem}
\newcommand{\omnameProblems}{Problems}
\newcommand{\omnameProof}{Proof}

% Admitted results and solutions appendix headers
\newcommand{\omadmittedtext}{Admitted at this level.}
\newcommand{\omsolutionof}{Solution of}
\newcommand{\ompage}{page}
\newcommand{\omnameGotoSolution}{Solution p.}

% Misc text macros
\newcommand{\st}{\text{ such that }}

% Cover / title page
\newcommand{\bookmaintitle}{One Computer Science Book}
\newcommand{\booktagline}{The One Computer Science Book to Rule Them All}
\newcommand{\bookdescription}{Computer science from first grade to the end of the bachelor's degree}
\newcommand{\bookcontributors}{The One Computer Science Book contributors}
\newcommand{\bookversionlabel}{Version}

% Colophon
\newcommand{\omcolophonsource}{Source code, issues and contributions:}
\newcommand{\omcolophonlicense}{Except where otherwise noted, this book is licensed under Creative Commons CC BY-NC-SA 4.0:}
\newcommand{\omcolophoncredit}{Please credit: One Computer Science Book, One Course (one-course.com).}

% Aside boxes (styles/onecs.sty): the recall box that opens a chapter
% building on another one, the history box, the dated box (its title reads
% "Facts as of 2026"), the safety box, the in-practice box.
\newcommand{\omnameRecall}{You already know}
\newcommand{\omnameHistory}{History}
\newcommand{\omnameDated}{Facts as of}
\newcommand{\omnameSafety}{Safe use}
\newcommand{\omnameInPractice}{In practice}
% Listings and captured outputs (\omcode, \omoutput).
\newcommand{\omnameListing}{Listing}
\newcommand{\omnameListings}{Listings}
\newcommand{\omnameOutput}{Output}
\newcommand{\omnameOutputs}{Outputs}
\newcommand{\omnameSomeLinesOut}{Some lines are left out.}

% Solutions appendix part title
\newcommand{\omsolutionspart}{Solutions to the Exercises}

% Part titles (Grades 1–12; used by shared part.tex files)
\@namedef{omstr@part.grade1}{Grade 1 (Age 6--7)}
\@namedef{omstr@part.grade2}{Grade 2 (Age 7--8)}
\@namedef{omstr@part.grade3}{Grade 3 (Age 8--9)}
\@namedef{omstr@part.grade4}{Grade 4 (Age 9--10)}
\@namedef{omstr@part.grade5}{Grade 5 (Age 10--11)}
\@namedef{omstr@part.grade6}{Grade 6 (Age 11--12)}
\@namedef{omstr@part.grade7}{Grade 7 (Age 12--13)}
\@namedef{omstr@part.grade8}{Grade 8 (Age 13--14)}
\@namedef{omstr@part.grade9}{Grade 9 (Age 14--15)}
\@namedef{omstr@part.grade10}{Grade 10 (Age 15--16)}
\@namedef{omstr@part.grade11}{Grade 11 (Age 16--17)}
\@namedef{omstr@part.grade12}{Grade 12 (Age 17--18)}

% Part titles (Bachelor Years 1–3; used by parts/bachelor-N/part.tex)
\@namedef{omstr@part.bachelor1}{Bachelor Year 1 (Age 18--19)}
\@namedef{omstr@part.bachelor2}{Bachelor Year 2 (Age 19--20)}
\@namedef{omstr@part.bachelor3}{Bachelor Year 3 (Age 20--21)}

% Parts inside the university volumes (parts/bachelor-N/part.tex).
\@namedef{omstr@part.b1.c}{C and the Machine Model}
\@namedef{omstr@part.b1.tools}{Tools}
\@namedef{omstr@part.b1.reasoning}{Reasoning about Programs}
\@namedef{omstr@part.b1.structures}{Data Structures}
\@namedef{omstr@part.b1.design}{Algorithm Design}
\@namedef{omstr@part.b1.machine}{Machine and Data}
\@namedef{omstr@part.b2.cpp}{C++}
\@namedef{omstr@part.b2.rust}{Rust and Functional Programming}
\@namedef{omstr@part.b2.architecture}{Architecture}
\@namedef{omstr@part.b2.systems}{Systems Programming}
\@namedef{omstr@part.b2.networks}{Networks}
\@namedef{omstr@part.b2.languages}{Languages and Logic}
\@namedef{omstr@part.b2.algorithms}{Algorithms II and Data}
\@namedef{omstr@part.b3.compilers}{Compilers and Programming Languages}
\@namedef{omstr@part.b3.os}{Operating Systems}
\@namedef{omstr@part.b3.theory}{Theory of Computation}
\@namedef{omstr@part.b3.data}{Data and Distributed Systems}
\@namedef{omstr@part.b3.security}{Security}
\@namedef{omstr@part.b3.building}{Building Software}
% Book 1's appendix of typing drills.
\@namedef{omstr@part.typing}{Typing Drills}

% Plural names + \omnameChapter, used by the \crefname declarations in
% onecs.sty. Without these, cleveref falls back to its English built-ins and
% a French book says "Theorem 3.1" / "Chapter 5" in the middle of French prose.
\newcommand{\omnameDefinitions}{Definitions}
\newcommand{\omnameTheorems}{Theorems}
\newcommand{\omnamePropositions}{Propositions}
\newcommand{\omnameLemmas}{Lemmas}
\newcommand{\omnameCorollarys}{Corollaries}
\newcommand{\omnameMethods}{Methods}
\newcommand{\omnameExamples}{Examples}
\newcommand{\omnameNotations}{Notations}
\newcommand{\omnameRemarks}{Remarks}
\newcommand{\omnameChapter}{Chapter}
\newcommand{\omnameChapters}{Chapters}
% Section, for \cref{sec:…} (2026-10-04: it had no \crefname, so cleveref
% printed its English built-in "Section" in every edition).
\newcommand{\omnameSection}{Section}
\newcommand{\omnameSections}{Sections}

% LaTeX's own structural names. babel would normally localise these, but it is
% optional here (see onecs.sty) — so without these a French book prints
% "Chapter 1" above every chapter title and "Appendix A" over the solutions.
\renewcommand{\chaptername}{Chapter}
\renewcommand{\partname}{Part}
\renewcommand{\contentsname}{Contents}
\renewcommand{\appendixname}{Appendix}
\renewcommand{\indexname}{Index}

